% Complete proofs for the measurable probability foundation in P15-U1.
\section{Measurable probability: proofs}
\label{sec:annex-p15-measurable-probability}

These proofs use the definitions in
Chapter~\ref{sec:p15-measurable-probability}. Measures may have infinite
values unless a finite or probability measure is specified. Every series
below has nonnegative terms and uses the supremum definition of that
chapter. The source foundations are Durrett's author version of
11 January 2019 \cite[secs.~1.1--1.3; app.~A.1]{DurrettPTEAuthor2019}.

\begin{annexitem}{annex:p15.generated-events}{thm:p15-generated-events}{Closure and the generated sigma-algebra}
\begin{proof}
Since \(\Omega\in\mathcal F\), complementation gives
\(\varnothing\in\mathcal F\). If \(A_n\in\mathcal F\), then
\[
 \bigcap_n A_n=\left(\bigcup_n A_n^c\right)^c\in\mathcal F.
\]
This also gives the prescribed empty intersection. For \(A,B\in\mathcal F\),
\(A\setminus B=A\cap B^c\) is measurable.

Let \((\mathcal H_i)_{i\in I}\) be a nonempty family of sigma-algebras
on \(\Omega\), with no countability restriction on \(I\). Each
\(\mathcal H_i\) contains \(\Omega\), so their intersection does.
If \(A\) belongs to every \(\mathcal H_i\), then \(A^c\) does.
If every \(A_n\) belongs to every \(\mathcal H_i\), closure in each
\(\mathcal H_i\) puts \(\bigcup_n A_n\) in every one of them.
This proves the intersection is a sigma-algebra. Applied to the defining
family for \(\sigma(\mathcal C)\), it proves existence. Every member
of that family contains \(\mathcal C\), so their intersection does.
By definition it is contained in each of them. Any other smallest
sigma-algebra with that property both contains and is contained in this
one, proving uniqueness.
\end{proof}
\end{annexitem}

\begin{annexitem}{annex:p15.measure-laws}{thm:p15-measure-laws}{Consequences of countable additivity}
\begin{proof}
For finitely many disjoint sets, append empty sets to obtain a sequence
and apply countable additivity. The added terms are zero. For an empty
family the assertion is \(\mu(\varnothing)=0\). If \(A\subseteq B\),
the sets \(A\) and \(B\setminus A\) are disjoint with union \(B\).
Their measures add to \(\mu(B)\). Nonnegativity proves monotonicity.
When \(\mu(B)<\infty\), monotonicity makes both summands finite,
so ordinary real subtraction gives the difference formula.

For an arbitrary sequence \((A_n)\), put \(D_1=A_1\) and
\(D_n=A_n\setminus\bigcup_{j<n}A_j\) for \(n>1\).
Each \(D_n\) is measurable and contained in \(A_n\). The \(D_n\)
are pairwise disjoint: membership in \(D_n\) excludes every earlier
\(A_j\). Their union equals \(\bigcup_n A_n\), since every point
in the latter union has a least index at which it occurs. Therefore
\[
 \mu\left(\bigcup_n A_n\right)=\sum_n\mu(D_n)
       \leq\sum_n\mu(A_n).
\]
The inequality follows first for each finite initial sum and then for
their suprema. If all \(A_n\) are null, the right side is zero.

Suppose now that \(A_n\) increases to \(A\). Set
\(D_1=A_1\) and \(D_n=A_n\setminus A_{n-1}\) for \(n>1\).
The \(D_n\) are disjoint, their union is \(A\), and their first
\(n\) members have union \(A_n\). Hence
\[
 \mu(A)=\sum_{j\geq1}\mu(D_j)
       =\sup_n\sum_{j=1}^n\mu(D_j)=\sup_n\mu(A_n).
\]
Monotonicity makes the last sequence increasing, so its supremum is its
extended real limit.

If instead \(A_n\) decreases to \(A\), the sets
\(C_n=A_1\setminus A_n\) increase to \(A_1\setminus A\).
With \(\mu(A_1)<\infty\), the difference formula gives finite
quantities throughout and
\[
 \mu(A_n)=\mu(A_1)-\mu(C_n)
     \longrightarrow\mu(A_1)-\mu(A_1\setminus A)=\mu(A).
\]
Monotonicity gives the decreasing direction. No infinite subtraction
has been performed.

For a probability measure, monotonicity with \(A\subseteq\Omega\)
gives the upper bound one, and the disjoint decomposition
\(\Omega=A\cup A^c\) gives the complement formula. If all \(A_n\)
have probability one, all \(A_n^c\) are null; their union is null, so
the complement \(\bigcap_n A_n\) has probability one. Finally,
\(\mathbb P(A_1)\leq1<\infty\), supplying the hypothesis needed
for continuity from above in every probability space.
\end{proof}
\end{annexitem}

\begin{annexitem}{annex:p15.countable-atomic}{thm:p15-countable-atomic}{Grouping nonnegative terms and constructing atomic measures}
\begin{proof}
First note the following finite grouping fact. If \(J_1,\ldots,J_m\)
are disjoint index sets, each carrying nonnegative terms, then
\[
 \sum_{j=1}^m\sup_{K\subseteq J_j\text{ finite}}\sum_{i\in K}a_i
 =\sup_{\substack{K_j\subseteq J_j\text{ finite}\\1\leq j\leq m}}
       \sum_{j=1}^m\sum_{i\in K_j}a_i.
\]
Every expression on the right is bounded above by the left. If all
\(m\) suprema are finite and \(m>0\), choose for each one a finite
sum within \(\varepsilon/m\) of it; their total is within
\(\varepsilon\) of the left side. Letting \(\varepsilon>0\) tend
to zero proves the reverse inequality. If a supremum is infinite, its
finite sums either include an infinite term or exceed every prescribed
finite bound. Choosing that set and empty sets for the others shows
that the right side is infinite too. When \(m=0\), both sides are zero.

Write \(L=\sum_{i\in I}a_i\) and
\(b_n=\sum_{i\in I_n}a_i\). For a finite \(J\subseteq I\),
only finitely many blocks \(I_n\) meet \(J\), and
\[
 \sum_{i\in J}a_i
   =\sum_{n:J\cap I_n\ne\varnothing}\sum_{i\in J\cap I_n}a_i
   \leq\sum_{n\geq1}b_n.
\]
Taking the supremum over \(J\) proves \(L\leq\sum_n b_n\).
For the reverse inequality, fix a finite collection \(M\) of block
indices. The finite grouping fact expresses \(\sum_{n\in M}b_n\)
as a supremum of sums over finite subsets chosen within those blocks.
Their unions are finite subsets of \(I\); each resulting sum is at
most \(L\). Thus \(\sum_{n\in M}b_n\leq L\). Supremizing over
finite \(M\) proves \(\sum_n b_n\leq L\). This argument includes
empty index sets and infinite values.

For the proposed \(\mu_p\), the value at the empty set is zero.
Given disjoint \(A_n\subseteq\Omega\), apply the proved grouping
identity to \(I=\bigcup_n A_n\), with blocks \(I_n=A_n\) and
terms \(p(\omega)\). This gives countable additivity exactly.
Nonnegativity is part of the sum definition, so \(\mu_p\) is a
measure. If the total mass is one it is a probability measure.

Conversely, let \(\mathbb P\) be a probability measure on all subsets
of a countable \(\Omega\). Put \(p(\omega)=\mathbb P(\{\omega\})\).
These are finite nonnegative numbers. The singleton decomposition of
\(\Omega\), with empty terms appended if \(\Omega\) is finite,
shows that their sum is one. Decomposing any \(A\subseteq\Omega\)
into its singletons similarly gives
\(\mathbb P(A)=\sum_{\omega\in A}p(\omega)\).
The grouping result ensures this value is independent of any enumeration.
Any inducing mass function must have value \(\mathbb P(\{\omega\})\)
at \(\omega\), establishing uniqueness.
\end{proof}
\end{annexitem}

\begin{annexitem}{annex:p15.borel-generators}{thm:p15-borel-generators}{Generating Borel sets by rectangles and rays}
\begin{proof}
There are countably many rectangles with rational endpoints in fixed
finite dimension: their parameters form a subset of \(\mathbb Q^{2d}\),
a finite product of countable sets. Each such rectangle is open and
hence Borel. If \(U\) is open and \(x\in U\), choose a ball about
\(x\) contained in \(U\). Density of the rationals allows rational
endpoints on both sides of each coordinate of \(x\), close enough that
the resulting rectangle lies in that ball. Consequently \(U\) equals
the union of all rational rectangles contained in \(U\): one inclusion
is immediate, and the other follows from this choice for every \(x\).
This is a countable union because the whole collection of rational
rectangles is countable. Thus every open set belongs to the sigma-algebra
they generate. Its equality with the Borel sigma-algebra follows from
the two inclusions.

In dimension one, let \(\mathcal H\) be generated by the rational
rays \((-\infty,q]\). These rays are Borel, since their complements
are open, so \(\mathcal H\subseteq\mathcal B(\mathbb R)\).
For every real \(a\),
\[
 (-\infty,a)=\bigcup_{\substack{q\in\mathbb Q\\q<a}}(-\infty,q],
 \qquad
 (-\infty,a]=\bigcap_{\substack{q\in\mathbb Q\\q>a}}(-\infty,q].
\]
For the first identity, a real number smaller than \(a\) lies below
a rational number still smaller than \(a\). For the second, a real
number larger than \(a\) lies above some rational still larger than
\(a\), and is excluded by that ray. These observations also verify
the other inclusions in both identities. Thus both sorts of real ray
belong to \(\mathcal H\), and
\((a,b)=(-\infty,b)\setminus(-\infty,a]\) belongs to it.
The first part of the proof now gives
\(\mathcal B(\mathbb R)\subseteq\mathcal H\).
The collection of all real closed rays contains the rational ones and
consists entirely of Borel sets, so it generates the same sigma-algebra.
\end{proof}
\end{annexitem}

\begin{annexitem}{annex:p15.measurability-criterion}{thm:p15-measurability-criterion}{Inverse images and the generated information}
\begin{proof}
For any function \(X:\Omega\to S\), inverse images satisfy
\[
 X^{-1}(S)=\Omega,\qquad
 X^{-1}(S\setminus B)=\Omega\setminus X^{-1}(B),\qquad
 X^{-1}\left(\bigcup_n B_n\right)=\bigcup_n X^{-1}(B_n).
\]
These follow by checking membership of an arbitrary \(\omega\).
For example, \(X(\omega)\) lies in a union exactly when it lies in
one of its members. Therefore
\(\mathcal D=\{B\subseteq S:X^{-1}(B)\in\mathcal F\}\)
is a sigma-algebra on \(S\). If the inverse images of all generators
are measurable, \(\mathcal C\subseteq\mathcal D\), whence
\(\sigma(\mathcal C)\subseteq\mathcal D\) by minimality. This
is measurability of \(X\). The converse follows because generators
are members of the target sigma-algebra. The Borel generator theorem
proves the asserted rational-ray criterion.

The same inverse-image identities show directly that
\(\{X^{-1}(B):B\in\mathcal S\}\) is a sigma-algebra on
\(\Omega\). It makes every target inverse image measurable by its
definition, and any sigma-algebra making \(X\) measurable must contain
all these inverse images. This proves its minimality and the inclusion
in \(\mathcal F\) when \(X\) is \(\mathcal F\)-measurable.

Finally set
\(\mathcal G=\sigma(\{X^{-1}(C):C\in\mathcal C\})\).
The inclusion \(\mathcal G\subseteq\sigma(X)\) follows because
every generator on the left belongs to the sigma-algebra on the right.
Apply the already proved criterion with \(\mathcal G\) as the
source sigma-algebra. It makes \(X\) measurable into
\((S,\sigma(\mathcal C))\), so minimality of \(\sigma(X)\)
gives the reverse inclusion.
\end{proof}
\end{annexitem}

\begin{annexitem}{annex:p15.measurable-operations}{thm:p15-measurable-operations}{Measurable operations and simple functions}
\begin{proof}
Let \(X:(\Omega,\mathcal F)\to(S,\mathcal S)\) and
\(f:(S,\mathcal S)\to(T,\mathcal T)\) be measurable. For
\(D\in\mathcal T\), \(f^{-1}(D)\in\mathcal S\), and
\[
 (f\circ X)^{-1}(D)=X^{-1}(f^{-1}(D))\in\mathcal F.
\]
Thus composition preserves measurability.

Suppose the coordinates \(X_1,\ldots,X_d\) are real random variables.
The inverse image of a rational open rectangle under their joint map is
\[
 \{(X_1,\ldots,X_d)\in\textstyle\prod_{j=1}^d(a_j,b_j)\}
       =\bigcap_{j=1}^d\{a_j<X_j<b_j\},
\]
a finite intersection of events. The generator theorem makes the joint
map measurable. Conversely the \(j\)-th coordinate projection has open
inverse images of open sets: near a vector whose \(j\)-th coordinate
is in an open set, a sufficiently small ball keeps that coordinate there.
It is therefore Borel measurable by the generator criterion. Composing
it with a measurable vector proves that its \(j\)-th coordinate is
measurable.

For a continuous map between Euclidean spaces, the inverse image of an
open set is open: continuity at a point of the inverse image provides a
small ball mapped into a ball contained in the target open set. Since
open sets generate the target Borel sigma-algebra, the criterion again
gives Borel measurability. Composition with a measurable vector gives
the continuous-function assertion. Sums and products are continuous by
the elementary real limit laws, absolute value satisfies
\(\bigl||x|-|y|\bigr|\leq|x-y|\), and
\[
 \max(x,y)=\frac{x+y+|x-y|}{2},\qquad
 \min(x,y)=\frac{x+y-|x-y|}{2}.
\]
Thus these particular functions are included.

For an indicator, the inverse image of any Borel set is one of
\(\varnothing,A,A^c,\Omega\), according to which of zero and one
that set contains. It is therefore measurable. A finite linear
combination of indicators is measurable by the finite operations just
proved and has at most finitely many values, since each indicator has
two possible values. The empty sum is constant and measurable.
Conversely, if a measurable function has finite range
\(\{c_1,\ldots,c_m\}\) with distinct \(c_j\), its fibers
\(A_j=\{X=c_j\}\) are measurable: singletons are Borel. They
partition \(\Omega\), and
\(X=\sum_{j=1}^m c_j\mathbf1_{A_j}\) pointwise. This proves the
equivalence of the simple-function descriptions.
\end{proof}
\end{annexitem}

\begin{annexitem}{annex:p15.measurable-limits}{thm:p15-measurable-limits}{Extended limits and increasing simple approximations}
\begin{proof}
For an extended measurable \(Z\), the strict sublevel events are
measurable because
\[
 \{Z<a\}=\bigcup_{\substack{q\in\mathbb Q\\q<a}}\{Z\leq q\}.
\]
This identity includes \(Z=-\infty\); \(Z=\infty\) belongs to
neither side. Complements give measurable events \(\{Z>a\}\) and
\(\{Z\geq a\}\). For a sequence of extended variables,
\[
 \{\sup_n Z_n\leq a\}=\bigcap_n\{Z_n\leq a\},\qquad
 \{\inf_n Z_n<a\}=\bigcup_n\{Z_n<a\}.
\]
The second equality follows because an infimum is strictly below \(a\)
exactly when at least one term is strictly below \(a\). Also
\[
 \{\inf_n Z_n\leq a\}
     =\bigcap_{k\geq1}\{\inf_n Z_n<a+1/k\}.
\]
These identities and the extended generator criterion prove measurability
of both the supremum and the infimum. Apply these conclusions first to
each tail and then to the sequence of tail values to prove measurability
of \(\liminf\) and \(\limsup\).

At any outcome where a real sequence has finite limit \(x\), for
every \(\varepsilon>0\) all sufficiently late terms lie between
\(x-\varepsilon\) and \(x+\varepsilon\). Their tail infima and
suprema therefore have common limit \(x\). Hence an everywhere finite
pointwise limit agrees with the measurable extended liminf. Its events
\(\{X\leq a\}\) are measurable, and the real generator criterion
makes it a real random variable.

For real \(X_n,X\), each \(\{|X_n-X|<1/k\}\) is measurable by
the preceding finite-operation theorem. The displayed countable
intersection and unions are therefore measurable. At an outcome in the
displayed event, every \(1/k\) eventually bounds the distance. Given
\(\varepsilon>0\), choose \(k\) with \(1/k<\varepsilon\) to
obtain convergence. Conversely, convergence supplies an eventual strict
bound \(1/k\) for every fixed \(k\). This proves the event identity.

For the approximation, every interval event and the terminal event
\(\{Z\geq n\}\) in the definition of \(s_n\) is measurable by
the sublevel and complement facts above. They form a finite partition
because \(Z\geq0\); the terminal event includes \(Z=\infty\).
Thus \(s_n\) is nonnegative, simple, and finite valued. Fix an outcome
and write its value as \(z\in[0,\infty]\). If \(z\geq n\),
then \(s_n=n\). If \(z<n\), there is exactly one integer
\(0\leq k<n2^n\) with \(k2^{-n}\leq z<(k+1)2^{-n}\), and
\(s_n=k2^{-n}\). In particular \(s_n\leq z\).

To check monotonicity, if \(z<n\), the grid of width \(2^{-(n+1)}\)
refines the grid of width \(2^{-n}\), retaining its lower endpoints.
The lower endpoint selected at stage \(n+1\) is therefore at least
the one selected at stage \(n\). If \(n\leq z<n+1\), the new grid
has the integer \(n\) as an endpoint, so its chosen value is at least
\(n=s_n\). If \(z\geq n+1\), the values are respectively \(n\)
and \(n+1\). These cases include \(z=\infty\), proving
\(s_n\leq s_{n+1}\) everywhere. For finite \(z\), every integer
\(n>z\) satisfies \(0\leq z-s_n<2^{-n}\), so \(s_n\to z\).
For \(z=\infty\), \(s_n=n\to\infty\). This proves the full
approximation assertion without any integration theorem.
\end{proof}
\end{annexitem}

\begin{annexitem}{annex:p15.pushforward-laws}{thm:p15-pushforward-laws}{The law as a measure on values}
\begin{proof}
Measurability of \(X\) makes \(\mathbb P_X\) defined on every
\(B\in\mathcal S\). Its values are nonnegative,
\(X^{-1}(\varnothing)=\varnothing\), and \(X^{-1}(S)=\Omega\),
so its empty-set value is zero and its total value is one. If \(B_n\)
are pairwise disjoint, then their inverse images are pairwise disjoint:
one value \(X(\omega)\) cannot lie in two different members. The
inverse-image union identity and countable additivity of \(\mathbb P\)
give
\[
 \mathbb P_X\left(\bigcup_n B_n\right)
   =\mathbb P\left(\bigcup_n X^{-1}(B_n)\right)
   =\sum_n\mathbb P(X^{-1}(B_n))=\sum_n\mathbb P_X(B_n).
\]
Thus \(\mathbb P_X\) is a probability measure.

For a measurable \(f\), the composite is measurable by the earlier
theorem, and for \(D\in\mathcal T\),
\[
 \mathbb P_{f(X)}(D)
   =\mathbb P(X^{-1}(f^{-1}(D)))=\mathbb P_X(f^{-1}(D)).
\]
In a countable atomic source space, substituting the defining atomic
measure into \(\mathbb P(X^{-1}(B))\) gives the stated outcome sum.
Finally, for countable \(S\) with all subsets measurable, its law is
a probability measure on that countable space. Apply the already proved
atomic representation theorem to \(\mathbb P_X\). Its singleton
weights are precisely \(p_X(x)\), giving both the event-sum formula
and normalization, including empty fibers and empty value events.
\end{proof}
\end{annexitem}

\begin{annexitem}{annex:p15.measure-uniqueness}{thm:p15-measure-uniqueness}{Pi-lambda argument and uniqueness}
\begin{proof}
An intersection of lambda-systems on \(S\) is a lambda-system:
each contains \(S\), and each of the two closure conditions persists
in the intersection. There is at least one lambda-system containing
\(\mathcal C\), namely \(2^S\). Let \(\mathcal L_0\) be the
intersection of all such lambda-systems. It is the smallest one.

A lambda-system \(\mathcal L\) closed under binary intersections is
a sigma-algebra. Indeed, \(S\setminus A\) belongs to it by nested
difference. Complementation and binary intersections give binary unions,
then finite unions by induction. For \(A_n\in\mathcal L\), the
finite unions \(\bigcup_{j=1}^n A_j\) increase, so their union belongs
to \(\mathcal L\) by its last defining property.

For \(A\in\mathcal L_0\), define
\(\mathcal G_A=\{B\subseteq S:A\cap B\in\mathcal L_0\}\).
It is a lambda-system. Its whole-set condition holds because
\(A\cap S=A\in\mathcal L_0\). If \(B\subseteq D\) belong to
\(\mathcal G_A\), then
\[
 A\cap(D\setminus B)=(A\cap D)\setminus(A\cap B)\in\mathcal L_0,
\]
since the latter two sets are nested members of \(\mathcal L_0\).
If \(B_n\) increases, then \(A\cap B_n\) increases and its union
is \(A\cap\bigcup_n B_n\), proving the increasing-union condition.

When \(A\in\mathcal C\), the pi-system property gives
\(\mathcal C\subseteq\mathcal G_A\), and minimality of
\(\mathcal L_0\) gives \(\mathcal L_0\subseteq\mathcal G_A\).
Thus \(A\cap B\in\mathcal L_0\) for \(A\in\mathcal C\) and
\(B\in\mathcal L_0\). By symmetry, for each fixed
\(A\in\mathcal L_0\), all \(B\in\mathcal C\) now belong to
\(\mathcal G_A\). Since this is a lambda-system, minimality again
gives \(\mathcal L_0\subseteq\mathcal G_A\). Hence
\(\mathcal L_0\) is closed under all binary intersections. By the
preceding paragraph it is a sigma-algebra, so
\(\sigma(\mathcal C)\subseteq\mathcal L_0\). Every lambda-system
containing \(\mathcal C\) contains \(\mathcal L_0\), proving the
first assertion. This also covers an empty pi-system.

Let \(\nu_1,\nu_2\) be the two probability measures in the second
assertion. The collection
\[
 \mathcal L=\{B\in\sigma(\mathcal C):\nu_1(B)=\nu_2(B)\}
\]
is a lambda-system on \(S\). Both measures give \(S\) value one.
For nested members \(A\subseteq B\), finite subtraction gives
\(\nu_i(B\setminus A)=\nu_i(B)-\nu_i(A)\), so their differences
have equal measure. For increasing \(B_n\) on which they agree,
continuity from below shows agreement on \(\bigcup_n B_n\).
The assumed agreement gives \(\mathcal C\subseteq\mathcal L\).
The first assertion yields \(\sigma(\mathcal C)\subseteq\mathcal L\),
which is equality of the measures on their entire domain.
\end{proof}
\end{annexitem}

\begin{annexitem}{annex:p15.cdf-laws}{thm:p15-cdf-laws}{Distribution limits, jumps and determination of the law}
\begin{proof}
For \(a\leq b\), \(\{X\leq a\}\subseteq\{X\leq b\}\);
monotonicity of probability proves that of \(F_X\). Its values lie
in \([0,1]\). Since \(X\) takes finite real values everywhere,
\(\{X\leq n\}\) increases to \(\Omega\) and
\(\{X\leq-n\}\) decreases to \(\varnothing\). The two
continuity results give \(F_X(n)\to1\) and \(F_X(-n)\to0\).
Monotonicity then gives the respective limits for unrestricted real
arguments tending to positive and negative infinity.

For fixed real \(a\), the events \(\{X\leq a+1/n\}\) decrease
to \(\{X\leq a\}\), so \(F_X(a+1/n)\to F_X(a)\).
If \(a<t<a+1/n\), monotonicity places \(F_X(t)\) between
\(F_X(a)\) and \(F_X(a+1/n)\). This proves right continuity.
The events \(\{X\leq a-1/n\}\) increase to \(\{X<a\}\),
so their probabilities increase to \(\mathbb P(X<a)\). Their
supremum is \(\sup_{t<a}F_X(t)\): every \(t<a\) is at most
\(a-1/n\) for sufficiently large \(n\). Thus
\(F_X(a-)=\mathbb P(X<a)\).

The disjoint decomposition
\(\{X\leq b\}=\{X\leq a\}\cup\{a<X\leq b\}\)
gives the interval formula by finite subtraction. Similarly,
\(\{X\leq a\}=\{X<a\}\cup\{X=a\}\) gives the jump formula.
For completeness, a bounded nondecreasing \(F\) has the left limit
specified by its supremum: for any \(\varepsilon>0\), choose
\(t<a\) with \(F(t)>F(a-)-\varepsilon\); every \(s\) between
\(t\) and \(a\) has its value between that lower bound and
\(F(a-)\). Together with right continuity, this proves that \(F_X\)
is continuous at \(a\) exactly when \(F_X(a-)=F_X(a)\).
By the jump formula this is exactly the condition \(p_X(a)=0\).

For \(k\geq1\), let \(D_k=\{a\in\mathbb R:p_X(a)\geq1/k\}\).
It has at most \(k\) elements. Otherwise choosing \(k+1\) distinct
members would give disjoint singleton events with total probability at
least \((k+1)/k>1\). Every positive mass is at least \(1/k\) for
some positive integer \(k\). Therefore the set of positive masses
is \(\bigcup_{k\geq1}D_k\), a countable union of finite sets, and
is at most countable. This counts all discontinuities by the preceding
characterization.

Equal laws give equal values on all rays, hence equal cumulative
functions. Conversely, if two cumulative functions agree, their laws
agree on the pi-system \(\{(-\infty,a]:a\in\mathbb R\}\): the
intersection of two such rays has the smaller endpoint and is another
member. This pi-system generates \(\mathcal B(\mathbb R)\) by the
Borel generator theorem. The probability measure uniqueness theorem
therefore makes the two laws equal on all Borel sets.
\end{proof}
\end{annexitem}
